% Generated draft: requires mathematical and source review.
\documentclass[11pt]{article}
\usepackage[margin=25mm]{geometry}
\usepackage{fontspec}
\setmainfont{Latin Modern Roman}
\usepackage{amsmath,amssymb,mathtools,unicode-math}
\setmathfont{Latin Modern Math}
\usepackage{enumitem,needspace,xurl,hyperref}
\setlength{\parindent}{0pt}
\setlength{\parskip}{5pt}
\setlength{\emergencystretch}{3em}
\newcommand{\N}{\mathbb{N}}
\newcommand{\Z}{\mathbb{Z}}
\newcommand{\Q}{\mathbb{Q}}
\newcommand{\R}{\mathbb{R}}
\newcommand{\C}{\mathbb{C}}
\newcommand{\F}{\mathbb{F}}
\newcommand{\A}{\mathbb{A}}
\newcommand{\PP}{\mathbb{P}}
\newcommand{\E}{\mathbb{E}}
\newcommand{\eps}{\varepsilon}
\newcommand{\dd}{\,\mathrm{d}}
\begin{document}
\section*{UnsolvedMath ID: 9700001}
\textbf{Source problem code:} catalogue entry
\subsection{Definitions and mathematical statement}
Let \((\Omega,\mathcal F,\Pr,(\mathcal F_i)_{i=0}^n)\) be a filtered probability space, and let \(X_0,\ldots,X_n\) be an adapted integrable real process. A stopping time \(T\) takes values in \(\{0,\ldots,n\}\) with \(\{T\leq i\}\in\mathcal F_i\). A martingale satisfies \(\E[X_{i+1}\mid\mathcal F_i]=X_i\) almost surely, and finite-horizon optional stopping gives \(\E X_T=\E X_0\) for every such \(T\).

Aldous asks for a mathematically useful definition of being a ``martingale for practical purposes'' with the following operational implication: if a process fails the definition, it should be practical to find a stopping time \(T\) such that
\[
 \E X_T\ne\E X_0.
\]
The inequality may have either sign; no positivity, payoff threshold, or tolerance is stipulated. The aim is effective discovery of a stopping rule, not merely abstract existence.

For a discrete process \(X_0,\ldots,X_n\), the related question is whether a natural polynomial-size collection of stopping-time constraints, of the form \(\E X_T=\E X_0\), can define a tractable class with this practical interpretation. The underlying information filtration, encoding or access model for the process law, size parameter for the polynomial, computational resources and required numerical accuracy are not specified. A polynomial number of constraints does not by itself ensure that their expectations can be computed or a violated constraint found efficiently. Likewise a nonzero but arbitrarily tiny bias need not be practically detectable; the source does not provide a quantitative replacement for its exact inequality.

This is a definitional and algorithmic programme. It does not prescribe an existing approximate-martingale class, a finite list of tests, or an error bound to adopt.
\subsection{Short English statement}
Define a useful practical analogue of a martingale whose failure enables a feasible stopping rule to change expected value. For a finite process, ask whether polynomially many natural stopping-time tests can yield such a tractable class.
\subsection{Sources}
\begin{itemize}
\item Local UnsolvedMath registry, permanent ID 9700001, source code catalogue entry; attributed to David Aldous; no proposal year supplied; update 2026-08-06.
\item Aldous, \emph{Open problems}, item fields.html: \url{https://www.stat.berkeley.edu/~aldous/Research/OP/index.html}. Registry access date 2026-07-29; HTML text with linked primary-PDF text. The background reports 23 active pages, with explicit delegated subproblems separated and compound bullets retained. Present primary-literature status and redistribution terms need review.
\item The supplied summary mentions quasi-martingales and approximate optional-stopping criteria in finance only as neighboring questions, with no exact reference or accepted definition solving this problem. Its unresolved-status assessment is unverified.
\end{itemize}
\end{document}
